\chapter{Discussion}
\label{chap:Discussion}

The evaluation showed that the prototype could run the planned participant sessions, collect complete module-level data, and expose several differences between tasks, input modes, and perturbation implementations. The results in Chapter~\ref{chap:results} were not uniform across all conditions, but they were also not uniformly negative. The controller condition, where tracking and confirmation were most stable, showed perturbation shifts above the participant-specific no-effect reference in every tested task/effect combination. The purpose of this discussion is therefore to interpret what the evaluation reveals about the framework: which parts of the prototype supported measurable prism-adaptation-like shifts, which variables made interpretation harder, and which implementation dimensions should be made explicit in future VR prism-effect studies.

The most important finding is methodological. The system gathered enough information to show that the observed baseline-to-post shifts were not determined by the perturbation type alone. The same effect mode could appear different depending on whether the participant used controllers or embodied hand tracking, whether the task was Line Bisection or Open-Loop Pointing, and whether the participant completed exposure carefully or quickly. This supports the project's central premise: a VR prism-adaptation system cannot be described adequately by saying only that a ``prism effect'' was applied. The perturbation class, input representation, confirmation method, visibility schedule, assessment task, and procedure all shape what the data mean.

\section{Interpreting the Noisy Results}
\label{sec:discussion-noisy-results}

The Results chapter showed measurable baseline-to-post perceived-centre shifts across the dataset, but the no-effect condition was also often non-zero. This makes the results difficult to interpret as a simple perturbation-versus-control comparison. If a participant changes their response centre in the None condition, then a later perturbation condition has to be compared against that participant's own ordinary drift rather than against zero. This is why the normalized values in Table~\ref{tab:results-condition-summary}, Figure~\ref{fig:results-effect-minus-none-linebisection}, and Figure~\ref{fig:results-effect-minus-none-openloop} are more informative than the absolute shifts alone.

The clearest condition-level pattern was found in the controller conditions. In Line Bisection, Translation, Rotation, and Skew all had positive median shifts above the participant-specific None condition, with Rotation showing the largest median difference. The same above-None direction was also present for controller Open-Loop Pointing, although the differences were smaller. This is the strongest evidence in the evaluation that the product worked as a prism-effect prototype: when the input mode tracked reliably, the perturbation modules produced larger baseline-to-post shifts than the same participants' no-effect modules. This should still be read as exploratory evidence because of the small sample size, but it is a clear positive pattern rather than a null result.

In contrast, the embodied conditions were harder to interpret because the None condition itself was often large and because embodied input introduced additional variability. This does not mean that the embodied implementation produced no adaptation. It means that, in this sample, any adaptation signal was harder to separate from tracking instability, hand posture, and no-effect drift.

The sample size also limits the strength of the conclusions. Ten participants produced 80 completed modules, which is a substantial amount of logged runtime data for a prototype evaluation, but it is still a small participant sample. With only five participants in each input mode, one rushed participant, one calibration problem, or one unusually stable participant can visibly change the group-level distribution. The boxplots and participant-level plots in Figure~\ref{fig:results-absolute-shifts}, Figure~\ref{fig:results-effect-minus-none-linebisection}, Figure~\ref{fig:results-effect-minus-none-openloop}, Figure~\ref{fig:results-participants-1-2}, and Figure~\ref{fig:results-participants-3-4} should therefore be read as exploratory evaluation data rather than as final statistical evidence.

\section{Participant Behaviour as a Source of Variability}
\label{sec:discussion-participant-behaviour}

Several participant behaviours affected the clarity of the data. Some participants completed the exposure phase exceedingly quickly and missed many targets. Participant~2 is the clearest example, as shown in Table~\ref{tab:results-participant-quality} and Figure~\ref{fig:results-quality-examples}. A fast pace is not automatically invalid, but in this setup it meant that the participant could move through exposure with many misses and limited stable correction. This reduces confidence that the exposure phase induced the intended sensorimotor recalibration.

Focus and motivation also varied. Two participants were observed as both fast and bored, and they began talking during the task. This matters because the task depends on repeated, consistent aiming behaviour. When participants divide attention between the task and conversation, the measured endpoint may reflect attention and compliance as much as adaptation. Fatigue was another visible factor. Some participants became tired over the repeated modules, especially because the protocol required many raised-arm actions. Fatigue could affect pointing precision, posture, and willingness to maintain the instructed movement style.

Body compensation was another important behavioural variable. Some participants used the upper body, wrist swivel, or shoulder movement to aim, even after being instructed not to. In one case, upper-body use appeared to make the participant very precise. This is problematic for interpretation because the prototype measured where the ray hit, not whether the participant used the intended reaching strategy. If a participant can compensate by rotating the wrist or torso, then the measured endpoint can look stable even if the intended prism-adaptation movement is not being performed.

The evaluation also showed posture differences in the embodied condition. Some participants held the hand sideways or pointed upward when instructed to point forward. In embodied hand tracking, small differences in hand posture could produce large differences in the projected aim direction. One participant repeatedly aimed upward while trying to point forward, and several temporary corrections were needed before the aiming direction became usable. This suggests that embodied hand input is not only a different visual representation, but also a different motor-control problem for the participant.

\section{Simulator Sickness and Session Tolerance}
\label{sec:discussion-ssq}

The SSQ results in Table~\ref{tab:results-ssq-summary} suggest that simulator sickness was present but not the dominant limitation of the evaluation. The mean total SSQ score increased after the session, and some participants reported more nausea, oculomotor discomfort, or disorientation after testing. However, the free-text comments and facilitator observations more often pointed to arm fatigue, aiming consistency, calibration, task pace, and tracking than to sickness symptoms alone.

This matters because session tolerance was still a procedural constraint even when simulator sickness was moderate. The evaluation required repeated raised-arm pointing across eight modules, so discomfort could appear as physical fatigue or reduced focus rather than only as classical cybersickness. The SSQ therefore supports a cautious interpretation: the prototype did not appear unusable because of simulator sickness, but future evaluations should still monitor sickness and fatigue separately because both can affect pointing behaviour and participant compliance.

\section{Embodied Tracking and Input Robustness}
\label{sec:discussion-embodied-tracking}

The embodied condition introduced technical and behavioural variability that was less visible in the controller condition. Some participants had intermittent tracking issues, while others completed embodied modules without major tracking problems. The issue was therefore not that embodied tracking always failed, but that it added a condition-specific source of instability. This matters because the prototype used a raycast from the tracked hand representation: small changes in finger angle, palm orientation, or tracking confidence could produce comparatively large changes in the projected aim. This matches prior virtual prism-adaptation work using Leap Motion hand tracking, where occlusion and self-occlusion were noted as practical problems and controllers were suggested as a possible way to reduce missing-hand errors at the cost of other trade-offs \cite{cho2022feasibility}.

The embodied condition should also not be interpreted as simply ``more prism-like'' because it shows a hand. Gerken et al. found that virtual hand feedback could speed early adaptation compared with cursor feedback, but did not produce corresponding differences in aftereffects, generalization, or intermanual transfer \cite{gerken2025sense}. This supports the interpretation that embodiment may affect the adaptation process and user experience without guaranteeing a larger post-exposure shift in every assessment task. In the present evaluation, the more stable controller condition produced the clearer above-None pattern, while the embodied condition exposed the practical tracking and posture problems that would need to be solved before it could be interpreted with the same confidence.

The final prototype used opposite-hand trigger confirmation instead of dwell confirmation. This was practically important because a two-second dwell across hundreds of attempts would have made the evaluation much slower and more tiring. The audio click feedback added to confirmation was well received by participants, because it made accepted responses clearer without requiring the participant to look away from the task. However, confirmation method remains a methodological variable. Several VR studies use controller buttons, dwell, physical touch, foot pedals, or inferred endpoints, and these are not equivalent interaction designs.

The controller implementation was more robust during testing, but it also has its own limitation. A controller trigger or button press can disturb the endpoint at the moment of measurement. This was one reason the prototype explored opposite-hand confirmation and other mitigation options. The literature is mixed here: several VR prism-adaptation studies use controller-button endpoint registration in some form \cite{Bourgeois2021,Gammeri2020,anan2025comparative}, while other systems avoid endpoint button presses through physical reaching, dwell, haptics, or inferred movement endpoints \cite{cho2022feasibility,gerken2025sense,Ramos2019}. This supports treating action-confirmation method as part of the framework rather than as a minor implementation detail.

\section{Reaching Movement and Ray-Based Aiming}
\label{sec:discussion-reaching-ray}

One major difference between the prototype and conventional physical prism-adaptation procedures concerns the movement itself. In many conventional setups, participants point from the body midline toward a target. Serino et al. describe physical pointing with the right index finger from the sternum \cite{serino2006mechanisms}, and Ladavas et al. compare terminal and concurrent prism adaptation using a wooden-box setup with physical pointing movements \cite{Serino2011}. In those procedures, the participant's arm movement is the action being adapted.

The prototype instead implemented a pointer raycast from the virtual controller or hand representation. This made the tasks feasible in VR and allowed the same target wall to support Open-Loop Pointing, Line Bisection, and Exposure. However, it also meant that participants could aim by changing wrist angle, controller orientation, arm posture, or torso posture. The endpoint was therefore not only a consequence of a forward reaching movement. Even though participants were instructed to keep the elbow bent and avoid excessive wrist swivel, the system did not enforce this mechanically.

This distinction is important for interpreting the results. If the participant solves the task by rotating the wrist or upper body, then the perturbation may be corrected through a strategy that differs from the reaching movement used in traditional prism-adaptation exposure. The prototype still measures a baseline-to-post spatial shift, but the motor process underlying that shift may not be the same as in classical PA. Future versions should either constrain the action more closely to a reaching movement or explicitly log and analyze wrist, arm, and torso contributions as separate variables.

\section{Exposure Feedback and Miss Handling}
\label{sec:discussion-exposure-feedback}

The literature review indicates that exposure visibility is not a trivial detail. Ladavas et al. directly compared terminal prism adaptation and concurrent prism adaptation and found that the two procedures had different effects in patients with neglect \cite{Serino2011}. Cho et al. implemented a virtual equivalent of terminal exposure by making the virtual hand visible only in a target-visible area \cite{cho2022feasibility}. Carter et al. similarly hid the initial portion of the reaching movement in their virtual-shift procedure because visual feedback during reaching can reduce prism-adaptation induction \cite{Carter2016}. These examples show that when the participant sees the hand or endpoint during exposure is part of the experimental procedure, not merely a visual-design choice.

Our implementation used concurrent visual feedback during exposure. Participants could see the perturbed hand or controller representation while aiming, and this was especially important for making the Skew effect understandable after the final implementation fixes. This was a stability-oriented design choice: the evaluation prioritized a consistent visible-feedback implementation across all effect modes rather than introducing an additional terminal-visibility rule that would itself become another uncontrolled implementation variable. However, based on the terminal-versus-concurrent comparison in Ladavas et al. \cite{Serino2011}, this is a meaningful procedural difference from classical terminal prism adaptation. The difference may matter most for Skew, because Skew is closest to a full-field or representation-level visual manipulation and may depend strongly on when the participant sees the displaced hand or controller.

The handling of misses is another procedural difference worth noting. In the prototype, a missed exposure attempt was still counted as an attempt, and the system advanced to a new target. This avoided forcing participants with poor accuracy to repeat indefinitely, which would have increased fatigue. The literature does not show one universal rule across all VR and physical protocols. Classical and VR protocols generally structure exposure as a sequence of pointing movements or trials, but the exact handling of missed or inaccurate trials is often not described in enough detail to know whether a target is repeated until hit. Anan et al. used 96 exposure reaches, while Carter et al. and Bourgeois et al. used 100 exposure reaches in their virtual-shift or virtual-prism procedures \cite{anan2025comparative,Carter2016,Bourgeois2021}. Our 60-attempt exposure blocks were therefore shorter than several prior protocols, and they should be interpreted as 60 attempts rather than 60 successful corrections. Hit rate, shown in Table~\ref{tab:results-participant-quality} and Figure~\ref{fig:results-quality-examples}, is therefore essential for interpreting exposure quality.

\section{Assessment Feedback}
\label{sec:discussion-assessment-feedback}

The pre- and post-exposure assessment tasks were intended to measure the participant's perceived middle or pointing endpoint without giving corrective feedback that would allow immediate error correction. This design choice is broadly consistent with the role of open-loop and post-exposure tests in the reviewed literature. Serino et al. distinguish exposure error reduction from post-exposure after-effect \cite{serino2006mechanisms}, and several VR studies include no-feedback or reduced-feedback test blocks for measuring aftereffects \cite{Gammeri2020,gerken2025sense}.

The literature review did not identify terminal corrective error feedback during baseline or post-exposure Open-Loop Pointing and Line Bisection as a required feature. Instead, corrective visual feedback is mainly an exposure-phase variable. Not providing terminal error feedback in the baseline and post blocks should therefore not be treated as a mistake in itself. The more important limitation is that the prototype still used a visible aiming ray or response representation in some tasks, which may have acted as a feedforward aiming aid. This is different from giving explicit error correction after a response, but it still affects how ``open-loop'' the task truly is.

\section{Task Transfer and Carry-Over}
\label{sec:discussion-transfer-carryover}

Several participant profiles suggested possible transfer between modules. For example, one module's post-exposure state appeared to bleed into the next module's baseline, which could make a later None condition appear shifted. This is visible in the results because the None condition was not consistently near zero. Carry-over is especially plausible in a protocol where all eight modules are completed back-to-back with only short breaks. Participant instruction and expectation may also matter in VR adaptation: Wähnert and Schäfer found that different instructions influenced the persistence of aftereffects under a virtual displacement, even when the physical task structure was otherwise comparable \cite{Wahnert2024}. This strengthens the need for highly standardized instructions between modules.

The literature often handles this problem with reset phases, washout periods, randomized or counterbalanced condition orders, or longer breaks. Ramos et al. used a reset step and breaks between conditions \cite{Ramos2019}. Anan et al. used a longer washout in a within-subject comparison \cite{anan2025comparative}. Our protocol prioritized feasibility and participant tolerance, so the full session was shorter and more continuous. The resulting data are useful for evaluating the prototype workflow, but carry-over makes it difficult to isolate each module as an independent adaptation effect.

The evaluation also tested each participant only once. This is sufficient for comparing how the framework behaves across tasks, effects, and input modes in one session, but it does not address long-term retention or repeated-session change. Those questions belong to later studies that use the framework in a rehabilitation-oriented protocol.

Future evaluations should treat carry-over as a first-class design issue. One option is to use a stronger washout or reset task between modules. Another is to increase the break time between perturbations. A third is to reduce each participant to fewer conditions, so that fewer perturbations are stacked in one session. This would reduce within-subject comparability, but it may produce cleaner interpretation of each effect.

\section{Product and Procedure Implications}
\label{sec:discussion-product-procedure}

The project aimed to facilitate prism-adaptation research and rehabilitation design through a more unified framework. Conceptually, the prototype moved toward that goal by making perturbation class, input mode, task, logging, and analysis explicit. The RShiny analysis pipeline was especially useful because it allowed the same session to be inspected from multiple levels: condition summaries, participant profiles, exposure quality, spatial response distributions, and individual trajectories. This made it possible to identify noisy modules and participant-specific issues rather than only reporting one endpoint number.

At the same time, the evaluation showed that the product alone is not enough. A more thorough procedure is also needed. Human factors such as posture, boredom, fatigue, calibration behaviour, and instruction compliance had large effects on the quality of the data. Some of these can be addressed through software, for example by adding posture checks, clearer calibration guidance, stricter pause logic, or warnings when hit rate becomes too low. Others require a better testing protocol, such as standardized verbal instructions, stronger familiarization, clearer exclusion criteria, or a fixed rule for when a module should be repeated.

This means that future work should improve both the prototype and the evaluation design. Improving only the visual perturbations would not solve the problems observed in the data if participants can still rush, compensate with the torso, or carry adaptation into the next module. Conversely, improving only the procedure would not solve implementation issues such as embodied tracking instability or unclear visibility schedules. The most useful next version would combine stricter procedural control with a software interface that logs the relevant control variables explicitly.

\section{Framework-Level Reflection}
\label{sec:discussion-framework-reflection}

The evaluation reinforces the need for a shared terminology framework in VR prism adaptation. Across the literature, different studies describe full-field shifts, effector-feedback shifts, virtual hand shifts, visuomotor rotations, and symbolic-body shifts using overlapping language \cite{Ramos2019,Gammeri2020,Bourgeois2021,cho2022feasibility,gerken2025sense}. Our own results show why this matters. Translation, Rotation, and Skew were all implemented as prism-like perturbations, but they did not produce identical patterns across tasks and input modes. The same label would hide those differences.

The prototype therefore contributes most clearly as a framework-building tool rather than as a finished rehabilitation product. It demonstrates that the relevant variables can be made explicit in code, logged per session, and visualized in a shared analysis dashboard. It also shows that several variables need to be treated as core descriptors rather than secondary details: perturbation class, exposure visibility, embodiment representation, action-confirmation method, task family, reset procedure, and participant posture.

The noisy data should not be read only as a weakness. The noise revealed which parts of the workflow are under-controlled and which variables need to be reported in future VR prism-adaptation studies. In that sense, the evaluation supports the broader project aim: before VR prism adaptation can become a robust rehabilitation framework, the field needs clearer language for what is being shifted, what the participant can see, how responses are confirmed, and how after-effects are measured.
